\chapter{Reference: \code{when} and conditionals}
\label{ch:refman-when}

\begin{aside}
A tiny helper for branching inside an element tree without
breaking the expression chain.
\end{aside}

\section{Forms}

\begin{itemize}[leftmargin=*]
  \item \code{when(cond, then\_element)} — render
    \code{then\_element} if \code{cond}, otherwise nothing.
  \item \code{when(cond, then, else)} — ternary form.
\end{itemize}

\section{Working example}

\begin{lstlisting}[language=C++]
#include <maya/maya.hpp>

using namespace maya;
using namespace maya::dsl;

struct WhenDemo {
    struct Model { bool on = false; int count = 0; };
    struct Toggle {};
    struct Bump {};
    struct Quit {};
    using Msg = std::variant<Toggle, Bump, Quit>;

    static Model init() { return {}; }

    static auto update(Model m, Msg msg)
        -> std::pair<Model, Cmd<Msg>>
    {
        return std::visit(overload{
            [&](Toggle) { m.on = !m.on; return std::pair{m, Cmd<Msg>{}}; },
            [&](Bump)   { m.count++;    return std::pair{m, Cmd<Msg>{}}; },
            [](Quit)    { return std::pair{Model{}, Cmd<Msg>::quit()}; },
        }, msg);
    }

    static Element view(const Model& m) {
        return v(
            when(m.on,
                t<"ON">  | Bold | Fg<100, 255, 100>,
                t<"OFF"> | Bold | Fg<255, 100, 100>),
            blank_,
            when(m.count > 0,
                text(std::format("clicks: {}", m.count))),
            when(m.count >= 5,
                t<"high five!"> | Bold | Fg<255, 220, 80>),
            blank_,
            t<"t toggle, b bump, q quit"> | Dim
        ) | pad<1> | border_<Round>;
    }

    static auto subscribe(const Model&) -> Sub<Msg> {
        return key_map<Msg>({
            {'q', Quit{}},
            {'t', Toggle{}},
            {'b', Bump{}},
        });
    }
};

int main() { run<WhenDemo>({.title = "when"}); }
\end{lstlisting}

\section{Notes}

\begin{itemize}[leftmargin=*]
  \item \code{when(false, e)} contributes nothing to layout.
  \item Both branches are \emph{constructed} either way;
    pick the cheap form if construction is expensive.
  \item For long branches, use a helper function and inline
    the call.
\end{itemize}

\section{Recap}

\begin{recap}
\begin{itemize}[leftmargin=*]
  \item \code{when} for inline branches.
  \item Two-arg "show if", three-arg "ternary".
  \item Both branches construct; only one renders.
\end{itemize}
\end{recap}

\section*{Workshop}

\begin{enumerate}[leftmargin=*]
  \item Use \code{when} to hide a footer when the window is
    too short.
  \item Render different views per "tab" using a sequence of
    \code{when}s.
  \item Replace a long \code{if/else} chain with nested
    \code{when}s.
\end{enumerate}
